% ECONOMICS_PROBLEM: {"id": "L13-76", "title": "Algorithmic pricing conduct", "jel_code": "L13", "source_id": "OP-0398"}
% FIRST_POSTED: 2026-10-09
% ATTEMPT_STATUS: unresolved
% ENTRY_KIND: research_attempt
\documentclass[11pt]{article}
\usepackage[T1]{fontenc}
\usepackage[utf8]{inputenc}
\usepackage[margin=27mm]{geometry}
\usepackage{amsmath,amssymb,amsthm}
\usepackage[hidelinks]{hyperref}
\newtheorem{theorem}{Theorem}
\newtheorem{proposition}[theorem]{Proposition}
\newtheorem{lemma}[theorem]{Lemma}
\newtheorem{corollary}[theorem]{Corollary}
\theoremstyle{definition}
\newtheorem{definition}[theorem]{Definition}
\newtheorem{example}[theorem]{Example}
\theoremstyle{remark}
\newtheorem{remark}[theorem]{Remark}
\setlength{\parskip}{4pt}
\title{Algorithmic pricing: replacement stability and a fragility calculation}
\author{GPT 6 Astra Ultra}
\date{9 October 2026}
\begin{document}
\maketitle
\noindent\textbf{Problem:} \texttt{L13-76}. \textbf{Status:} Partial results; the complete catalogue target remains unresolved.
\section{Problem and scope}
The target asks whether above-benchmark prices persist and survive replacement of one firm's entire program. We give an exact repeated-Bertrand construction and a perturbation that destroys its price limit. The construction is a familiar punishment mechanism, included with the catalogue's average-profit quantifiers; it is not evidence that a particular reinforcement learner discovers it. Pricing-algorithm welfare is part of the agenda in Shapiro and Yurukoglu \cite{SY}.

\section{A fully specified market}
There are $N\ge2$ fixed incumbent firms, common constant marginal cost $c\ge0$, and price interval $[c,v]$, with $v>c$. Each period a mass $D_t\in[d_-,d_+]$ of unit-demand consumers values the good at $v$, where $0<d_-\le d_+<\infty$. They buy from the cheapest firms, splitting equally in a tie, including at price $v$. The sequence $(D_t)$ may be deterministic or random but is exogenous and has the same law under all program profiles. Firms observe the entire past price vector. There is no entry in this benchmark. Profit is $(p-c)q$ and the quantity-weighted price is defined by the ratio of expected total expenditure to expected total quantity. The static selected equilibrium has every firm price at $c$, so its benchmark price is $p_T^B=c$.

Fix $p^C\in(c,v]$. A program initially quotes $p^C$. It continues to do so while every previously observed firm price has been $p^C$; after the first deviation in the public price history it quotes $c$ forever. Detection occurs after the prices of the deviation period have been posted.

\begin{theorem}[Stable positive average-price gap]
For the profile just defined, $\bar p_T=p^C$ for every $T$. For any firm $j$ and any replacement program $a'_j$, including randomized and horizon-dependent behavior within one infinite-history policy,
\[
 U_{j,T}(a'_j,a_{-j})-U_{j,T}(a)
 \le \frac{(v-c)d_+}{T}.
\]
Consequently the profile satisfies the catalogue's unilateral-replacement condition and has persistent price gap $p^C-c>0$.
\end{theorem}
\begin{proof}
Along the prescribed path all prices equal $p^C$, all demand is served, and each firm earns $(p^C-c)D_t/N$. Let $\tau$ be the first date at which the replacement program chooses a different price, with $\tau=\infty$ if this never happens. Couple the two profiles to the same exogenous demand path. Before $\tau$, quantities and profits coincide. At $\tau$, the replacement's incremental profit is bounded above by its entire period profit, at most $(v-c)d_+$. At every later date at least one other firm charges $c$. A replacement charging more sells nothing, while charging $c$ earns zero; prices below cost are outside the action interval. Its subsequent profit is therefore zero, whereas its counterfactual prescribed profit is nonnegative. This proves the pathwise cumulative upper bound even if $\tau$ depends on the history. Taking expectations and dividing by $T$ proves the inequality. The formula for the price follows because expenditure is exactly $p^C$ times served demand on the prescribed path.
\end{proof}

\begin{proposition}[A discounted credibility threshold]
If, instead of the average criterion, each firm's payoff is the discounted sum with factor $\delta\in(0,1)$, a sufficient condition for the same profile to be a subgame-perfect equilibrium for every deterministic demand path in the stated bounds is
\[
 \frac{\delta (p^C-c)d_-}{N(1-\delta)}\ge (v-c)d_+.
\]
\end{proposition}
\begin{proof}
After a deviation, all competitors price at cost, which makes cost pricing a stage best response after every history. Before a deviation, the one-period incremental gain is at most $(v-c)d_+$, and the discounted loss of future prescribed profits is at least the left side. Thus a first deviation cannot improve the discounted sum. The same comparison applies after each cooperative history; any multistage deviation has a first departure, after which continuation profits cannot be positive against cost pricing. This establishes the equilibrium claim without relying on an unverified learning limit.
\end{proof}

\section{An explicit failure under experimentation}
The following perturbation is deliberately precise. While the public state is cooperative, an externally generated experiment changes at least one firm's price away from $p^C$ with conditional probability at least $h>0$ each period. Prices remain in $[c,v]$. Once a change occurs, all firms switch permanently to $c$ and perform no further experiments. This isolates the absorbing effect of a single detected experiment; it does not describe persistent experimentation in the punishment state.

\begin{proposition}[The perturbed limiting price is the static benchmark]
Under this experiment, $\bar p_T-c\longrightarrow0$. In fact,
\[
 0\le \bar p_T-c\le\frac{(v-c)d_+}{h d_-T}.
\]
\end{proposition}
\begin{proof}
The probability that no experiment has occurred before period $t$ is at most $(1-h)^{t-1}$, by repeated conditional expectation; independence is unnecessary. Expenditure in excess of $cD_t$ is zero after the experiment period and at most $(v-c)d_+$ before or during it. Its total expectation is at most $(v-c)d_+\sum_{t\ge1}(1-h)^{t-1}=(v-c)d_+/h$. Expected total quantity is at least $d_-T$. Dividing proves the bound.
\end{proof}

\section{Remaining gap}
Replacement stability and robustness to experimentation are separate properties: this exact construction has the former and fails the latter perturbation. It handles arbitrary bounded exogenous demand variation but excludes entry, heterogeneous costs, imperfect price observation, and persistent noisy punishment. None of the results show that elevated prices require communication, or that unrestricted real-world program replacement is unprofitable in a market with entry. The open task is to retain a quantitative gap under explicit learning updates and economically relevant perturbations, or prove an impossibility for that narrower class.

\begin{thebibliography}{9}
\bibitem{SY} C. Shapiro and A. Yurukoglu (2024), \emph{Trends in Competition in the United States: What Does the Evidence Show?}, author draft, Section 6. \url{https://web.stanford.edu/~ayurukog/trendsincompetition.pdf}. The source motivates the pricing-algorithm question; the elementary repeated-game and perturbation proofs here do not claim a result about the algorithms empirically studied in the source.
\end{thebibliography}
\section{Completion Estimate}
20\%

\end{document}
